% section: 特殊関数の広がり

\section{広がり}


\subsection*{読み方}

漸化式を解くと,しばしば次のような疑問に出会う.

\begin{mdquote}
  この数列の答えが,毎回同じ種類の関数として現れるのはなぜだろうか.
\end{mdquote}

例えば,三角関数は加法定理を持つだけでなく,隣り合う整数倍の値について
\[
  \cos((n+1)\theta)
  =2\cos\theta\cos(n\theta)-\cos((n-1)\theta)
\]
という漸化式を満たす.
この漸化式を,\,$\theta$ を固定して $n$ の数列だと見れば,
三角関数は漸化式の解として現れている.

同じように,階乗を実数へ拡張しようとするとガンマ関数が現れ,
円や球面の波を記述しようとするとベッセル関数やルジャンドル多項式が現れる.
これらは,単に名前のついた公式の寄せ集めではない.

自然な問題を解こうとすると,
\begin{enumerate}
  \item まず漸化式や微分方程式が現れ,
  \item その規則に従う標準的な解を選び,
  \item その解が別の問題にも繰り返し現れるため,
  \item 一つの数学的対象として名前を持つようになる.
\end{enumerate}

この一連の流れを,特殊関数と呼ぶことができる.
ここでの目的は,特殊関数を暗記することではない.
「なぜその関数を作る必要があったのか」
「漸化式・微分方程式・積分・直交性が,同じ対象の別の顔としてどう現れるのか」
を追跡することである.

\subsection{翻訳}

第2章第7節では,倍角公式や対称式を使って数列を解いた.
ここでは,その解法がなぜ何度も現れるのかを,一つの共通原理として見直す.

\subsubsection{共役な表現}

漸化式を
\[
  a_{n+1}=F(a_n)
\]
と書く.直接 $F$ を繰り返し作用させる代わりに,
\[
  a_n=\varphi(t_n)
\]
と表して,より簡単な規則
\[
  t_{n+1}=G(t_n)
\]
を解けないか考える.

この置換が有効になる条件は,
\[
  F(\varphi(t))=\varphi(G(t))
  \tag{A.1}
\]
である.先に $\varphi$ で表してから $F$ を作用させても,
先に $G$ を作用させてから $\varphi$ で数に戻しても結果が同じになる.

\begin{mdquote}
\[
  a_1=\varphi(t_1),
  \qquad
  a_n=\varphi\left(G^{\,n-1}(t_1)\right)
\]
と予想したら,最後に初項と漸化式を帰納法で確認する.特殊関数型の
解法は,この「表現を変えて規則を簡単にする」操作として統一できる.
\end{mdquote}

\subsubsection{倍角}

例えば,
\[
  F(x)=2x^2-1,
  \qquad
  \varphi(t)=\cos t,
  \qquad
  G(t)=2t
\]
なら,倍角公式から (A.1) が成り立つ.したがって,
\[
  a_{n+1}=2a_n^2-1,
  \qquad
  a_1=\cos\theta
\]
の解は
\[
  a_n=\cos(2^{n-1}\theta)
\]
である.

初項が $1$ の範囲を外れると,同じ恒等式を持つ $\cosh$ に切り替える.
\[
  \cosh 2t=2\cosh^2t-1
\]
なので,\,$a_1=\cosh t$ なら
\[
  a_n=\cosh(2^{n-1}t)
\]
となる.これは別の公式を暗記する話ではなく,初項に合う表現を選ぶ話である.

\subsubsection{対称式}

三角関数を使わずに,
\[
  (u+u^{-1})^2-2=u^2+u^{-2}
  \tag{A.2}
\]
という恒等式を使うこともできる.\,$a_1=u+u^{-1}$ で
\[
  a_{n+1}=a_n^2-2
\]
なら,
\[
  a_n=u^{2^{n-1}}+u^{-2^{n-1}}
\]
である.指数の計算が簡単になり,元の数列の漸化式が解ける.

さらに,
\[
  (u+u^{-1})^3-3(u+u^{-1})=u^3+u^{-3}
\]
から,
\[
  a_{n+1}=a_n^3-3a_n
\]
には指数を $3$ 倍する表現が対応する.これらの多項式をまとめて扱う
視点が,後のチェビシェフ多項式につながる.

\subsubsection{近似}

置換によって一般項を得られない場合でも,漸化式を解へ近づくように設計できる.
$x^2=A$ を解くために接線を使うと,
\[
  a_{n+1}=\frac12\left(a_n+\frac{A}{a_n}\right)
  \tag{A.3}
\]
が得られる.これはニュートン法である.

ここでは一般項ではなく,誤差が次にどう変わるかを調べる.\,$r=\sqrt A$ と
おけば,
\[
  a_{n+1}-r=\frac{(a_n-r)^2}{2a_n}
\]
となり,十分近いところでは誤差が二乗されることが分かる.
漸化式は,既に決まった数列を読むものだけでなく,近似を設計する道具にもなる.

\subsection{構造}

\subsubsection{繰り返し現れる構造}

特殊関数には,すべてを一行で定める公式のような定義はない.
むしろ,次のような特徴をいくつも持つ関数を,特別に扱うのである.

\begin{itemize}
  \item 自然な漸化式を満たす.
  \item 微分方程式の標準的な解になる.
  \item 母関数や級数で生成できる.
  \item 積分表示を持つ.
  \item 直交性や対称性など,計算に使える構造を持つ.
  \item 多くの別分野の問題に現れる.
\end{itemize}

例えば,ある関数がたった一つの問題にしか現れないなら,
わざわざ名前をつける必要はない.
しかし,波,確率,近似,多項式,微分方程式の解に同じ関数が登場するなら,
その関数を研究すること自体に価値が生まれる.

\subsubsection{数列を関数へ広げる}

漸化式
\[
  a_{n+1}=F(n,a_n)
\]
は,整数 $n$ の上で定義された対象である.
ここで,\,$n$ を実数や複素数へ広げられないかと考える.

もちろん,すべての数列に自然な拡張があるわけではない.
しかし,拡張を要求すると,もとの漸化式だけでは決まらなかった
「どの関数を選ぶか」という問題が生じる.
この選択を,連続性,積分表示,成長条件,対称性などで決めていく.

特殊関数の物語は,
\[
  \text{整数上の漸化式}
  \longrightarrow
  \text{より広い変数を持つ関数}
\]
という拡張の物語でもある.

\subsection{ガンマ関数}

\subsubsection{階乗}

階乗は
\[
  n!=1\cdot2\cdot3\cdots n
\]
であり,
\[
  (n+1)!=(n+1)n!
\]
を満たす.
この漸化式だけを見ると,整数 $n$ に対して定義されている.

しかし,
\[
  \left(\frac12\right)!
\]
のような量を考えたい場面がある.
階乗を整数の外へ延長するには,どの関数を選べばよいだろうか.

\subsubsection{積分}

正の実数 $x$ に対して,
\[
  \Gamma(x)=\int_0^\infty t^{x-1}e^{-t}\,dt
\]
と定義する.
これをガンマ関数という.

なぜこの積分なのか.
まず,\,$t^{x-1}$ は指数 $x$ を自然に受け入れる.
また,\,$e^{-t}$ は $t\to\infty$ で急速に小さくなるため,
積分を収束させる.
さらに,部分積分をすると階乗の漸化式が現れる.

実際,\,$x>0$ に対して
\[
\begin{aligned}
  \Gamma(x+1)
  &=\int_0^\infty t^x e^{-t}\,dt\\
  &=\left[-t^xe^{-t}\right]_0^\infty
    +x\int_0^\infty t^{x-1}e^{-t}\,dt\\
  &=x\Gamma(x).
\end{aligned}
\]
\begin{mdquote}
  $\Gamma(x+1)=x\Gamma(x)$
\end{mdquote}
が得られる.

さらに,
\[
  \Gamma(1)=\int_0^\infty e^{-t}\,dt=1
\]
なので,正の整数 $n$ に対して帰納的に
\[
  \Gamma(n+1)=n!
\]
となる.
ガンマ関数は,階乗の漸化式を実数へ延長したものだと分かる.

\subsubsection{半整数}

ガンマ関数の面白さは,半整数の値を持つことである.
まず
\[
  \Gamma\left(\frac12\right)
  =\int_0^\infty t^{-1/2}e^{-t}\,dt
\]
において $t=u^2$ とおくと,
\[
  \Gamma\left(\frac12\right)
  =2\int_0^\infty e^{-u^2}\,du.
\]

ガウス積分
\[
  \int_{-\infty}^{\infty}e^{-u^2}\,du=\sqrt{\pi}
\]
を用いれば,
\[
  \Gamma\left(\frac12\right)=\sqrt{\pi}.
\]
したがって,
\[
  \Gamma\left(\frac32\right)
  =\frac12\Gamma\left(\frac12\right)
  =\frac{\sqrt{\pi}}2
\]
である.

ここで,階乗を連続化しただけのつもりが,
確率論や正規分布に現れる $\sqrt{\pi}$ とつながった.
特殊関数は,異なる分野を結ぶための共通言語になる.

\subsubsection{ベータ関数}

次の積分を考える.
\[
  B(x,y)=\int_0^1 t^{x-1}(1-t)^{y-1}\,dt
  \qquad (x,y>0).
\]
これをベータ関数という.

変数変換や二重積分の計算から,
\[
  B(x,y)=\frac{\Gamma(x)\Gamma(y)}{\Gamma(x+y)}
\]
が得られる.
この式は単なる公式ではない.
一つの積分に二つの指数が現れたことで,
ガンマ関数の積と,指数の和が同時に登場したのである.

階乗の積
\[
  \frac{m!n!}{(m+n)!}
\]
を連続化しようとするとベータ関数が現れる.
整数上の組合せ的な式を連続へ伸ばすと,
積分と特殊関数が現れるという一つの型が見える.

\subsection{超幾何関数}

\subsubsection{係数比}

数列 $a_n$ の隣り合う項の比
\[
  \frac{a_{n+1}}{a_n}
\]
を考える.
この比が $n$ の有理式になる数列を,
超幾何的な数列と呼ぶことがある.

例えば,
\[
  a_n=n!
\]
ならば
\[
  \frac{a_{n+1}}{a_n}=n+1.
\]
また,
\[
  a_n=\binom{m}{n}
\]
ならば
\[
  \frac{a_{n+1}}{a_n}
  =\frac{m-n}{n+1}.
\]

この見方では,数列の項を直接見るのではなく,
「一段進むと何倍されるか」を見る.
漸化式の最も基本的な視点が,特殊関数の級数へ拡張される.

\subsubsection{階乗記号}

次の記号を導入する.
\[
  (a)_n=a(a+1)(a+2)\cdots(a+n-1),
  \qquad (a)_0=1.
\]
これを上昇階乗,またはポッホハマー記号という.
一段の関係は
\[
  (a)_{n+1}=(a+n)(a)_n
\]
である.

一方,
\[
  a^{\underline n}
  =a(a-1)\cdots(a-n+1)
\]
を下降階乗という.
二項係数は
\[
  \binom{a}{n}
  =\frac{a^{\underline n}}{n!}
\]
と書ける.

通常のべき $a^n$ だけでなく,階乗型の積を基底として使うと,
組合せ論と級数を統一的に扱える.

\subsubsection{超幾何級数}

代表的な超幾何級数を
\[
  {}_2F_1(a,b;c;x)
  =
  \sum_{n=0}^{\infty}
  \frac{(a)_n(b)_n}{(c)_n\,n!}x^n
\]
と定義する.
係数を $u_n$ とすると,
\[
  \frac{u_{n+1}}{u_n}
  =
  \frac{(a+n)(b+n)}{(c+n)(n+1)}x
\]
であり,\,$n$ の有理式になっている.

この関数は,特殊関数を一つずつ別々に覚えるのではなく,
「係数比が有理式になる級数」という大きなクラスを作る.
多くの関数が,超幾何関数の特別な場合として書ける.

\subsection{チェビシェフへの橋}

チェビシェフ多項式の定義,漸化式,範囲による振る舞い,近似への応用は,
直前の第7節で詳しく扱った.ここでは,第2章第7節の解法との関係だけを確認する.

$x=\cos\theta$ とおけば,
\[
  T_n(x)=\cos(n\theta),
  \qquad
  T_{n+1}(x)=2xT_n(x)-T_{n-1}(x)
\]
である.つまり,本編で
\[
  a_{n+1}=2x a_n-a_{n-1}
\]
を解くために使った置換が,多項式の列そのものを定義している.

本編では $n$ に関する数列を解き,付録第7節では $x$ を変数とする
多項式族を調べた.同じ漸化式が,解法と理論という二つの入口を持っているのである.

\subsection{ルジャンドル多項式}

\subsubsection{三項漸化式}

ルジャンドル多項式 $P_n(x)$ は,
\[
  (n+1)P_{n+1}(x)
  =(2n+1)xP_n(x)-nP_{n-1}(x)
\]
を満たす.
\[
  P_0(x)=1,\qquad P_1(x)=x
\]
から,
\[
\begin{aligned}
  P_2(x)&=\frac12(3x^2-1),\\
  P_3(x)&=\frac12(5x^3-3x),\\
  P_4(x)&=\frac18(35x^4-30x^2+3)
\end{aligned}
\]
が得られる.

チェビシェフ多項式の漸化式とよく似ているが,係数が $n$ に依存している.
このわずかな違いが,直交性や微分方程式の違いを生む.

\subsubsection{母関数}

ルジャンドル多項式の母関数は
\[
  \frac1{\sqrt{1-2xt+t^2}}
  =\sum_{n=0}^{\infty}P_n(x)t^n
\]
である.

一つずつ $P_n$ を定義する代わりに,
全ての多項式を一つの関数に詰め込んだ.
左辺を $t$ について展開すれば,係数として全ての $P_n$ が取り出せる.

母関数を微分したり,\,$1-2xt+t^2$ を掛けたりして係数を比較すると,
三項漸化式が得られる.
つまり,母関数は漸化式を解くためだけでなく,
漸化式を生み出す装置でもある.

\subsubsection{微分方程式}

ルジャンドル多項式は,
\[
  (1-x^2)y''-2xy'+n(n+1)y=0
\]
の多項式解である.

同じ $P_n$ が,
\begin{itemize}
  \item 三項漸化式,
  \item 母関数,
  \item 二階微分方程式
\end{itemize}
という三つの方法で記述される.
どれが本当の定義で,どれが派生結果なのかは,目的によって変わる.
特殊関数を学ぶとは,同じ対象を異なる言語で読めるようになることである.

\subsubsection{直交性}

$[-1,1]$ 上の積分に対して,異なる次数のルジャンドル多項式は直交する.
\[
  \int_{-1}^{1}P_m(x)P_n(x)\,dx=0
  \qquad(m\neq n).
\]

直交性の意味は,図形的な「角度」と同じである.
ベクトルの成分を内積で取り出すように,
ある関数を $P_n$ の和として展開したとき,各係数を積分で取り出せる.

例えば,適切な関数 $f$ を
\[
  f(x)\simeq\sum_{n=0}^{N}c_nP_n(x)
\]
と近似するなら,
\[
  c_n=
  \frac{2n+1}{2}
  \int_{-1}^{1}f(x)P_n(x)\,dx
\]
となる.
多項式が単なる式の集まりから,
関数を分解するための座標系になった.

\subsection{エルミート多項式}

\subsubsection{定義}

ガウス関数
\[
  e^{-x^2}
\]
を何度も微分すると,
多項式とガウス関数の積が現れる.
そこで
\[
  H_n(x)e^{-x^2}
  =(-1)^n\frac{d^n}{dx^n}e^{-x^2}
\]
によってエルミート多項式 $H_n$ を定義する.

最初のものは
\[
\begin{aligned}
  H_0(x)&=1,\\
  H_1(x)&=2x,\\
  H_2(x)&=4x^2-2,\\
  H_3(x)&=8x^3-12x.
\end{aligned}
\]

\subsubsection{生成関数}

エルミート多項式の生成関数は
\[
  e^{2xt-t^2}
  =\sum_{n=0}^{\infty}H_n(x)\frac{t^n}{n!}
\]
である.
$t$ で微分すると,
\[
  (2x-2t)e^{2xt-t^2}
  =\sum_{n=0}^{\infty}H_{n+1}(x)\frac{t^n}{n!}.
\]
右辺の係数を比較すれば,
\[
  H_{n+1}(x)=2xH_n(x)-2nH_{n-1}(x)
\]
が得られる.

ここでも,生成関数は公式を圧縮するだけではない.
微分を一度行うことで,隣り合う多項式の関係を一気に作り出す.

\subsubsection{方程式と直交性}

エルミート多項式は
\[
  y''-2xy'+2ny=0
\]
を満たす.
また,重み $e^{-x^2}$ を用いると,
\[
  \int_{-\infty}^{\infty}
  H_m(x)H_n(x)e^{-x^2}\,dx=0
  \qquad(m\neq n).
\]

ルジャンドル多項式では区間 $[-1,1]$ と重み $1$ が現れた.
エルミート多項式では区間が全実数になり,重みが $e^{-x^2}$ になる.
どの区間で,どの重みを使うかが変わると,
それに適した直交多項式が生まれる.

\subsection{ベッセル関数}

\subsubsection{漸化式}

ベッセル関数 $J_\nu(x)$ は,
\[
  x^2y''+xy'+(x^2-\nu^2)y=0
\]
という微分方程式に現れる.
これは円筒座標で波動方程式を分離すると現れる方程式である.

整数次数については,隣り合う次数の間に
\[
  J_{n+1}(x)
  =\frac{2n}{x}J_n(x)-J_{n-1}(x)
\]
という漸化式がある.
係数に $1/x$ が含まれるため,チェビシェフ多項式のような多項式列ではない.
それでも,隣り合う解が一つの規則でつながっている.

\subsubsection{生成関数}

ベッセル関数は,形式的に
\[
  \exp\left(\frac{x}{2}\left(t-\frac1t\right)\right)
  =\sum_{n=-\infty}^{\infty}J_n(x)t^n
\]
という生成関数でまとめられる.

両辺を $t$ で微分し,左辺の微分を計算すると,
\[
  \frac{x}{2}\left(t+\frac1t\right)
  \sum_{n=-\infty}^{\infty}J_n(x)t^n
  =
  \sum_{n=-\infty}^{\infty}nJ_n(x)t^n
\]
となる.
係数を比較すると,
\[
  \frac{x}{2}\bigl(J_{n-1}(x)+J_{n+1}(x)\bigr)
  =nJ_n(x)
\]
が得られ,整理して
\[
  J_{n+1}(x)=\frac{2n}{x}J_n(x)-J_{n-1}(x)
\]
が戻る.

ここでは,漸化式を解いたのではなく,生成関数を設計して漸化式を取り出した.
特殊関数は,規則の解であるだけでなく,規則を生成する仕組みでもある.

\subsubsection{初期条件}

二階漸化式には,一般に二つの独立な解がある.
ベッセル方程式にも,原点付近で性質の異なる二つの解がある.
どちらを選ぶかは,次の条件で決まる.

\begin{itemize}
  \item 原点で有限か.
  \item 無限遠でどのように振る舞うか.
  \item 境界条件を満たすか.
  \item 対称性を持つか.
\end{itemize}

「漸化式を満たす」だけでは,解は一意に決まらない.
初期条件や境界条件を含めて,初めて自然な特殊関数が選ばれる.

\subsection{直交多項式と行列}

\subsubsection{構造}

チェビシェフ,ルジャンドル,エルミートなど,
多くの直交多項式は三項漸化式を満たす.
一般の形を
\[
  p_{n+1}(x)
  =(x-a_n)p_n(x)-b_np_{n-1}(x)
\]
と書く.

なぜ三項だけなのか.
直交性があると,\,$xp_n(x)$ を同じ直交多項式の基底で展開したとき,
多くの次数の項が消え,隣り合う三つの次数だけが残る.
三項漸化式は,直交性の代数的な影である.

\subsubsection{行列式}

三重対角行列
\[
  J_n=
  \begin{pmatrix}
    a_0&\sqrt{b_1}&&\\
    \sqrt{b_1}&a_1&\sqrt{b_2}&\\
    &\sqrt{b_2}&a_2&\ddots\\
    &&\ddots&\ddots
  \end{pmatrix}
\]
を考える.
$xI-J_n$ の行列式を $D_n(x)$ とすると,行列式の展開から
\[
  D_{n+1}(x)
  =(x-a_n)D_n(x)-b_nD_{n-1}(x)
\]
が得られる.

これは直交多項式の三項漸化式と同じ形である.
したがって,直交多項式の零点は,ある三重対角行列の固有値として現れる.

\subsubsection{零点}

多項式の零点を直接解くのは,高次数になるほど難しくなる.
しかし,その多項式が行列式として書けるなら,
零点は行列の固有値になる.

この見方の利点は多い.

\begin{itemize}
  \item 零点が実数になる理由を,対称行列の固有値から理解できる.
  \item 零点の範囲を,行列の評価で見積もれる.
  \item 数値計算で,根を直接探す代わりに固有値計算が使える.
\end{itemize}

漸化式から出発した多項式が,
行列のスペクトルという別の数学へ移る.
特殊関数を学ぶとき,関数の名前よりも,
どの構造が別の分野へ翻訳されるかを見る方が重要である.

\subsection{見取り図}

\subsubsection{四つの顔}

ここまでの例を,同じ関数に現れる構造として整理する.

\begin{center}
\begin{tabular}{c|c|c|c}
  対象 & 漸化式 & 生成の仕方 & 自然な問題\\ \hline
  $\Gamma$ & $\Gamma(x+1)=x\Gamma(x)$
    & 積分 & 階乗の連続化\\
  $T_n$ & $T_{n+1}=2xT_n-T_{n-1}$
    & 三角関数・生成関数 & 区間上の近似\\
  $P_n$ & 三項漸化式
    & 母関数 & 球面の対称性\\
  $H_n$ & $H_{n+1}=2xH_n-2nH_{n-1}$
    & ガウス関数の微分 & 正規分布・振動\\
  $J_n$ & $J_{n+1}=\frac{2n}{x}J_n-J_{n-1}$
    & 生成関数 & 円筒の波\\
\end{tabular}
\end{center}

この表の目的は暗記ではない.
「同じ関数が,問題によって別の顔を見せる」という事実を確認することである.

\subsubsection{定義}

チェビシェフ多項式は,
\[
  T_n(\cos\theta)=\cos(n\theta)
\]
で定義できる.
しかし,漸化式でも,複素数でも,微分方程式でも定義できる.

どの定義を採用するかは,何をしたいかで決まる.

\begin{itemize}
  \item 値を計算したいなら漸化式.
  \item 三角関数との関係を見たいなら定義式.
  \item 近似したいなら直交性.
  \item 理論をまとめたいなら微分方程式.
\end{itemize}

数学では,定義を一つに固定することよりも,
同じ対象へ複数の入口を持つことが強みになる.

\subsection{読み方}

\subsubsection{質問}

新しい漸化式に出会ったら,次の順に考えるとよい.

\begin{enumerate}
  \item 係数は $n$ に依存するか,変数 $x$ に依存するか.
  \item 次項との比は有理式になるか.
  \item 母関数を作ると,微分や積分で閉じるか.
  \item 三項漸化式なら,直交性や行列が隠れていないか.
  \item ある微分方程式の級数解として現れないか.
  \item 境界条件や対称性によって,自然な解が選ばれていないか.
\end{enumerate}

これらは,名前を当てるためのクイズではない.
どの構造を調べれば,漸化式の背後にある問題へ戻れるかを探す質問である.

\subsubsection{意味}

特殊関数は,初等関数より難しい関数を追加で覚える章ではない.
むしろ,初等関数だけでは表せない自然な現象に対して,
どのように標準的な解を選び,どの構造を保存し,どの計算を可能にするかを学ぶ章である.

漸化式から始めた読者にとって重要なのは,
\[
  \text{隣り合う項を結ぶ規則}
\]
が,
\[
  \text{関数の値を結ぶ規則}
\]
へ広がることだ.
そして,その関数が積分,微分方程式,行列,確率,波動の中で同じ姿を見せる.

\subsection*{まとめ}

この付録では,特殊関数を一覧として紹介するのではなく,
\[
  \text{階乗の拡張}
  \longrightarrow
  \text{三角関数の多項式化}
  \longrightarrow
  \text{直交性}
  \longrightarrow
  \text{波と行列}
\]
という流れを追った.各関数は,漸化式を解くために選んだ表現が,
積分,微分方程式,直交性,行列,波動へ広がった結果として現れている.
